%========================================================
% EXTRA APPENDIX — Basis Choice and Pauli Operator Representations
%========================================================

\section{Basis Choice and Pauli Operator Representations}

This appendix clarifies the distinction between abstract quantum operators and their matrix representations, with particular emphasis on the dependence of Pauli matrices on the chosen basis representation.

\medskip

A common source of confusion in quantum information theory is the apparent identification between Pauli operators and their standard matrix forms. In reality, the matrices commonly denoted by

\[
\sigma_x,
\qquad
\sigma_y,
\qquad
\sigma_z
\]

represent the coordinate expressions of abstract operators relative to a chosen basis.

The operators themselves are basis-independent linear transformations acting on the Hilbert space.

%--------------------------------------------------------
% Hilbert Space and Basis Choice
%--------------------------------------------------------

\subsection*{Hilbert Space and Basis Choice}

A single qubit is described by a two-dimensional complex Hilbert space

\[
\mathcal H \cong \mathbb C^2.
\]

To perform explicit calculations, one introduces an orthonormal basis.

In quantum information theory, the standard choice is the computational basis

\[
B_{cb}
=
\{
|0\rangle,
|1\rangle
\}.
\]

However, this basis is not fundamentally privileged by quantum mechanics. Any orthonormal pair of states may be chosen as the primary basis of the Hilbert space.

\medskip

For example, one may alternatively choose the basis

\begin{equation}
B_{+-}
=
\{
|+\rangle,
|-\rangle
\},
\label{eq:plus_minus_basis_appendix}
\end{equation}

where

\begin{equation}
|+\rangle
=
\frac{|0\rangle+|1\rangle}{\sqrt{2}},
\qquad
|-\rangle
=
\frac{|0\rangle-|1\rangle}{\sqrt{2}}.
\label{eq:plus_minus_states_appendix}
\end{equation}

The computational basis is therefore a convention rather than a physical necessity.

%--------------------------------------------------------
% Abstract Definition of Pauli Operators
%--------------------------------------------------------

\subsection*{Abstract Definition of Pauli Operators}

Pauli operators are fundamentally defined through their action on basis states.

For instance, in the computational basis, the Pauli \(X\) operator is defined by

\begin{equation}
\sigma_x |0\rangle = |1\rangle,
\qquad
\sigma_x |1\rangle = |0\rangle.
\label{eq:sigma_x_action_computational}
\end{equation}

Thus, \( \sigma_x \) acts as a basis-state exchange operator.

\medskip

Importantly, Eq.~\eqref{eq:sigma_x_action_computational} defines the operator itself, not its matrix representation.

The matrix representation appears only after coordinates are assigned to the basis vectors.

Using

\[
|0\rangle
=
\begin{pmatrix}
1\\
0
\end{pmatrix},
\qquad
|1\rangle
=
\begin{pmatrix}
0\\
1
\end{pmatrix},
\]

the matrix representation of \( \sigma_x \) in the computational basis becomes

\begin{equation}
[\sigma_x]_{B_{cb}}
=
\begin{pmatrix}
0 & 1 \\
1 & 0
\end{pmatrix}.
\label{eq:sigma_x_matrix_computational}
\end{equation}

%--------------------------------------------------------
% Basis Dependence of Matrix Representations
%--------------------------------------------------------

\subsection*{Basis Dependence of Matrix Representations}

The matrix representation of an operator depends on the chosen basis.

To illustrate this explicitly, consider the alternative basis \( B_{+-} \) introduced in Eq.~\eqref{eq:plus_minus_basis_appendix}.

Using Eq.~\eqref{eq:sigma_x_action_computational}, one obtains

\begin{align}
\sigma_x |+\rangle
&=
\sigma_x
\left(
\frac{|0\rangle+|1\rangle}{\sqrt{2}}
\right)
\nonumber
\\
&=
\frac{
\sigma_x|0\rangle
+
\sigma_x|1\rangle
}{\sqrt{2}}
\nonumber
\\
&=
\frac{
|1\rangle+|0\rangle
}{\sqrt{2}}
\nonumber
\\
&=
|+\rangle,
\end{align}

and similarly,

\begin{align}
\sigma_x |-\rangle
&=
\sigma_x
\left(
\frac{|0\rangle-|1\rangle}{\sqrt{2}}
\right)
\nonumber
\\
&=
\frac{
|1\rangle-|0\rangle
}{\sqrt{2}}
\nonumber
\\
&=
-|-\rangle.
\end{align}

Therefore, in the basis \( B_{+-} \),

\begin{equation}
[\sigma_x]_{B_{+-}}
=
\begin{pmatrix}
1 & 0 \\
0 & -1
\end{pmatrix}.
\label{eq:sigma_x_matrix_plus_minus}
\end{equation}

This matrix coincides with the standard matrix representation usually associated with \( \sigma_z \).

\medskip

Consequently, the same abstract operator admits different matrix representations depending on the chosen basis.

%--------------------------------------------------------
% Rotation of Bloch-Sphere Axes
%--------------------------------------------------------

\subsection*{Rotation of Bloch-Sphere Axes}

The previous result admits a direct geometric interpretation on the Bloch sphere.

When the computational basis \( \{|0\rangle,|1\rangle\} \) is chosen as the reference basis, the corresponding quantization axis defines the Bloch \(Z\)-axis.

The states \( |+\rangle \) and \( |-\rangle \) then lie along the Bloch \(X\)-axis.

\medskip

If the basis \( B_{+-} \) is instead chosen as the primary basis, the coordinate system of the Bloch sphere is effectively rotated:

\begin{itemize}

\item the former \(X\)-axis becomes the new \(Z\)-axis,

\item the former \(Z\)-axis becomes the new \(X\)-axis.

\end{itemize}

Consequently, the matrix representations associated with the Pauli operators are permuted.

\medskip

This demonstrates that the labels \(X\), \(Y\), and \(Z\) are not intrinsically tied to fixed matrices, but depend on the chosen reference basis.

%--------------------------------------------------------
% Pauli Operators in the B_{+-} Basis
%--------------------------------------------------------

\subsection*{Pauli Operators in the \(B_{+-}\) Basis}

Suppose now that \( B_{+-} \) is chosen as the primary computational basis of the Hilbert space.

One may then define new Pauli operators relative to this basis.

The operator exchanging the basis states is

\begin{equation}
X_{(+-)}
|+\rangle
=
|-\rangle,
\qquad
X_{(+-)}
|-\rangle
=
|+\rangle.
\label{eq:new_x_operator_plus_minus}
\end{equation}

Its matrix representation is therefore

\begin{equation}
[X_{(+-)}]_{B_{+-}}
=
\begin{pmatrix}
0 & 1 \\
1 & 0
\end{pmatrix}.
\label{eq:new_x_matrix_plus_minus}
\end{equation}

However, expressed in the original computational basis, this operator corresponds physically to the former \( \sigma_z \).

\medskip

Similarly, the diagonal operator in the \( B_{+-} \) basis is

\begin{equation}
Z_{(+-)}
|+\rangle
=
|+\rangle,
\qquad
Z_{(+-)}
|-\rangle
=
-|-\rangle,
\label{eq:new_z_operator_plus_minus}
\end{equation}

with matrix representation

\begin{equation}
[Z_{(+-)}]_{B_{+-}}
=
\begin{pmatrix}
1 & 0 \\
0 & -1
\end{pmatrix},
\label{eq:new_z_matrix_plus_minus}
\end{equation}

which corresponds physically to the former \( \sigma_x \).

\medskip

Thus, changing the primary basis rotates the correspondence between Bloch-sphere axes and Pauli operators while preserving the underlying qubit algebra.

%--------------------------------------------------------
% Physical Interpretation
%--------------------------------------------------------

\subsection*{Physical Interpretation}

In abstract quantum information theory, Pauli operators are initially defined at the purely algebraic level.

The physical interpretation arises only after associating the chosen basis with a specific experimental encoding.

For polarization qubits, one commonly identifies

\[
|0\rangle
\leftrightarrow
|H\rangle,
\qquad
|1\rangle
\leftrightarrow
|V\rangle,
\]

so that the computational basis corresponds to horizontal and vertical polarization states.

Only after this identification do the Bloch-sphere axes acquire direct physical meaning in terms of polarization measurements.

\medskip

Consequently, the matrix forms of the Pauli operators are not fundamentally physical objects. Rather, they represent coordinate descriptions of abstract operators relative to a chosen basis convention.

This distinction is essential for understanding why quantities such as concurrence remain invariant under local basis rotations despite the apparent presence of specific Pauli matrices in their formal definitions.

